% Part: propositional-logic
% Chapter: syntax-and-semantics
% Section: semantic-notions

\documentclass[../../../include/open-logic-section]{subfiles}

\begin{document}

\olfileid{pl}{syn}{sem}

\olsection{અર્થવિચારી ખ્યાલો}

આપણે નીચેના અર્થવિચારી ખ્યાલો વ્યાખ્યાયિત કરીએ છીએ:

\begin{defn}
\begin{enumerate}
\item કોઈ~$\pAssign{v}$ માટે $\pSat{v}{!A}$ હોય, તો
  !!^a{formula}~$!A$ \emph{સંતોષ્ય} છે; કોઈ પણ $\pAssign{v}$ માટે
  $\pSat{v}{!A}$ ન હોય, તો તે \emph{અસંતોષ્ય} છે;
\item બધાં !!{valuation}s~$\pAssign{v}$ માટે $\pSat{v}{!A}$ હોય, તો
  !!^a{formula}~$!A$ \emph{પુનરુક્તિ} છે;
\item !!^a{formula}~$!A$ સંતોષ્ય હોય પણ પુનરુક્તિ ન હોય, તો તે
  \emph{નિવાર્ય} છે;
\item જો $\Gamma$ એ !!{formula}sનો ગણ હોય, તો દરેક એવા
  !!{valuation}~$\pAssign{v}$ માટે કે જેના માટે $\pSat{v}{\Gamma}$,
  $\pSat{v}{!A}$ પણ હોય ત્યારે અને માત્ર ત્યારે $\Gamma\Entails !A$
  (``$\Gamma$, $!A$ને ફલિત કરે છે'').
\item જો $\Gamma$ એ !!{formula}sનો ગણ હોય, તો એવા કોઈ
  !!a{valuation}~$\pAssign{v}$ માટે $\pSat{v}{\Gamma}$ હોય ત્યારે
  $\Gamma$ \emph{સંતોષ્ય} છે; અન્યથા $\Gamma$ \emph{અસંતોષ્ય} છે.
\end{enumerate}
\end{defn}

\begin{prob}
 નીચેનાં ચાર !!{formula}sમાંનું દરેક (a)~સંતોષ્ય, (b)~પુનરુક્તિ અને
 (c)~નિવાર્ય છે કે નહિ તે નક્કી કરો.
\begin{enumerate}
  \item \( ( \Obj p_0 \lif ( \lnot \Obj p_1 \lif \lnot \Obj p_0 ) ) \).
  \item \( ( ( \Obj p_0 \land \lnot \Obj p_1 ) \lif ( \lnot \Obj p_0 \land \Obj p_2 )) \liff ( ( \Obj p_2 \lif \Obj p_0 ) \lif ( \Obj p_0 \lif \Obj p_1 )) \).
  \item \( ( \Obj p_0 \liff \Obj p_1 ) \lif ( \Obj p_2 \liff \lnot \Obj p_1 ) \).
  \item \( (( \Obj p_0 \liff ( \lnot \Obj p_1 \land \Obj p_2 )) \lor ( \Obj p_2 \lif ( \Obj p_0 \liff \Obj p_1 ))) \).
\end{enumerate}
\end{prob}

\begin{prop}
\ollabel{prop:semanticalfacts}
\begin{enumerate}
\item $!A$ પુનરુક્તિ હોય ત્યારે અને માત્ર ત્યારે
  $\emptyset \Entails !A$;
\item જો $\Gamma \Entails !A$ અને $\Gamma \Entails !A \lif !B$, તો
  $\Gamma \Entails !B$;
\item જો $\Gamma$ સંતોષ્ય હોય, તો $\Gamma$નો દરેક સાન્ત ઉપગણ પણ
  સંતોષ્ય છે;
\item \ollabel{def:monotonicity} એકદિશવર્ધિતા: જો
  $\Gamma \subseteq \Delta$ અને $\Gamma \Entails !A$, તો
  $\Delta \Entails !A$ પણ;
\item \ollabel{def:Cut} પરંપરિતતા: જો $\Gamma \Entails !A$ અને
  $\Delta \cup \{ !A\} \Entails !B$, તો $\Gamma \cup \Delta \Entails
  !B$.
\end{enumerate}
\end{prop}

\begin{proof}
અભ્યાસપ્રશ્ન.
\end{proof}

\begin{prob}
\olref[pl][syn][sem]{prop:semanticalfacts} સાબિત કરો.
\end{prob}

\begin{prop}\ollabel{prop:entails-unsat}
  $\Gamma \Entails !A$ ત્યારે અને માત્ર ત્યારે જ્યારે
  $\Gamma \cup \{\lnot !A\}$ અસંતોષ્ય હોય.
\end{prop}

\begin{proof}
અભ્યાસપ્રશ્ન.
\end{proof}

\begin{prob}
\olref[pl][syn][sem]{prop:entails-unsat} સાબિત કરો.
\end{prob}

\begin{thm}[અર્થાનુસારી નિગમન પ્રમેય]
  \ollabel{thm:sem-deduction} $\Gamma \Entails !A \lif !B$ ત્યારે અને માત્ર
  ત્યારે જ્યારે $\Gamma \cup \{!A\} \Entails !B$.
\end{thm}

\begin{proof}
અભ્યાસપ્રશ્ન.
\end{proof}

\begin{prob}
\olref[pl][syn][sem]{thm:sem-deduction} સાબિત કરો.
\end{prob}

\end{document}
